\subsection{Generalised CP Symmetry in Modular-Invariant Models of Flavour --- arXiv:1905.11970}

\textbf{Authors:} P. P. Novichkov, J. T. Penedo, S. T. Petcov, A. V. Titov

\paragraph{主な主張}
有限モジュラー対称性と一般化CPの整合条件を導き、モジュライ $\tau$ の真空期待値をCP破れの唯一の源とするフレーバー模型を構成できることを示した\cite{Novichkov:2019sqv}。

\paragraph{新規性}
モジュライと物質場へのCP作用をモジュラー変換と統一的に定式化し、$S_4$ レプトン模型の自由度と予言がCPによって制限されることを示した。

\paragraph{位置付け}
モジュラーフレーバー模型に一般化CPを課す際の基本的な定式化である。

\paragraph{コメント（読書メモ）}
読書メモ：内田さんに紹介していただいた論文。
